% year: תשפ"ד
% semester: ב
% moed: א
% number: 2
% categories: continuity, derivatives
% source: exams/מבחנים/תשפד/מועד א תשפד סמסטר ב.pdf

\begin{question}
נגדיר פונקציה באופן הבא:
\[ f(x)=\begin{cases} x^2\sin\left(\frac{3}{x}\right) & x\neq0\\ 0 & x=0\end{cases} \]
\begin{enumerate}
  \item האם $f(x)$ רציפה בנקודה $x=0$?
  \item האם $f(x)$ גזירה בנקודה $x=0$?
\end{enumerate}
\end{question}

\begin{hint}
$\sin\left(\frac3x\right)$ חסומה; השתמשו במשפט הסנדוויץ' (או "אפסה כפול חסומה").
\end{hint}

\begin{solution}
\begin{enumerate}
  \item לכל $x\neq0$ מתקיים $\left|\sin\frac3x\right|\le1$, ולכן
  \[ 0\le |f(x)|=x^2\left|\sin\frac3x\right|\le x^2. \]
  מכיוון ש-$x^2\to0$ כאשר $x\to0$, לפי משפט הסנדוויץ' $|f(x)|\to0$, כלומר $\lim_{x\to0}f(x)=0=f(0)$. לכן $f$ \textbf{רציפה} ב-$x=0$.
  \item נחשב את הנגזרת לפי ההגדרה:
  \[ f'(0)=\lim_{h\to0}\frac{f(h)-f(0)}{h}=\lim_{h\to0}\frac{h^2\sin\frac3h}{h}=\lim_{h\to0}h\sin\frac3h. \]
  שוב $0\le\left|h\sin\frac3h\right|\le|h|\to0$, ולפי משפט הסנדוויץ' הגבול קיים ושווה $0$ (פונקציה השואפת לאפס כפול פונקציה חסומה). לכן $f$ \textbf{גזירה} ב-$x=0$ ו-$f'(0)=0$.
\end{enumerate}
\end{solution}
